\section{Обзор литературы}
Среди методов управления самыми популярными являются Scrum и Kanban. Идея методологии Scrum - разработка итерациями. То есть за определенный промежуток времени команда должна выпустить часть готового продукта. Scrum традиционно применяют при создании нового продукта. Kanban, напротив, чаще используют в крупных компаниях с уже работающим продуктом и множеством команд. Этот метод сфокусирован на управлении потоками работ и ограничении незавершенных работ.

Принципы гибких методов управления сложно усвоить исключительно с помощью книг и статей. Эта проблема подтверждается результатами исследований в области менеджмента и бизнес-образования. Например, обзор \textit{The effectiveness of games for educational purposes: A review of recent research}~\cite{games-for-edu} подтверждает: игровые методы и симуляции развивают управленческие навыки лучше, чем лекции или чтение литературы. Это связано с необходимостью понимать внутренние процессы системы и их взаимосвязь.

В первую очередь это касается методологий Kanban~\cite{kanban-book-bor} и Lean~\cite{lean-book}. Поэтому в обучении широко применяются игровые симуляции. В игре участники могут тестировать любые решения и сразу видеть их последствия, не боясь навредить реальному проекту. Эффективность этого метода подтверждается теорией из книги Д. Колба \textit{Experiential Learning Cycle}~\cite{kolb}. Согласно его модели, глубокое обучение проходит через четыре этапа: накопление конкретного опыта, рефлексивное наблюдение, абстрактную концептуализацию и активное экспериментирование. Бизнес-игра позволяет быстро пройти этот цикл, закрепляя ценные практические навыки.

В работе \textit{Learning with a strategic management simulation game: A case study}~\cite{learning-with-simulation} авторы также изучали эффективность бизнес-симуляций. Исследователи делают вывод, что игры помогают увидеть картину целиком и использовать теорию в деле. Так проще понять связи внутри процесса, чем при изучении обычных учебных материалов.

Рассмотрим подробнее эти подходы. В своей книге Дэвид Андерсон определяет шесть базовых практик метода Kanban~\cite{kanban-book-and}:
\begin{enumerate}
    \item Визуализация потока задач.
    \item Ограничение объема незавершенной работы (WIP).
    \item Непосредственное управление потоком.
    \item Четкие правила выполнения работы.
    \item Регулярные встречи для сбора обратной связи.
    \item Совместное улучшение процесса.
\end{enumerate}
а также подчеркивает принцип: \textit{start with what you do now}. В игровой форме участникам наглядно видно связь принятых решений в управлении и скорость доведения задач до конца. Этот метод тесно связан с концепцией Lean. Целью обоих методов является устранение потерь и увеличение пропускной способности. Практическая часть опирается на Теорию ограничений Э. Голдратта~\cite{gold-goal-book}. Согласно ей, общая производительность системы лимитируется её "узкими местами". Эти принципы легли в основу разрабатываемой симуляции, моделирующей реальные процессы для обучения. 

Также в управлении проектами часто используют закон Литтла~\cite{little-law}. На практике этот закон позволяет менеджерам проектов точнее прогнозировать сроки и понимать реальную скорость работы команды. В методологии Kanban закон выражается формулой:
\begin{equation}
    \label{little-law-formula}
    L = \lambda\times W
\end{equation}
где $L$ - это незавершенная работа (количество активных задач в системе), $\lambda$ - пропускная способность (количество выполненных задач за единицу времени), $W$ - время цикла (время прохождения задачи через процесс). Однако, прочитав, сложно понять, как она работает на самом деле, людьми эта формула воспринимается абстрактно. Игровая симуляция дает людям понять на примере, как действует этот закон, и позволяет увидеть, какой результат дают ограничения на WIP. 

На данный момент есть онлайн-версия игры Featureban\footnote{Онлайн версия игры доступна по ссылке: \url{https://tools.flowfast.io/featureban}, дата обр. 13.01.2026}, но эта версия игры только для одной команды и рассчитана на 3–5 человек. Версия игры на несколько команд существует только в оффлайн формате и разработана сертифицированными экспертами из России (Neogenda)\cite{multi-featureban}. где $L$ - это незавершенная работа (количество активных задач в системе), $\lambda$ - пропускная способность (количество выполненных задач за единицу времени), $W$ - время цикла (время прохождения задачи через процесс). Однако, прочитав, сложно понять, как она работает на самом деле, людьми эта формула воспринимается абстрактно. Игровая симуляция дает людям понять на примере, как действует этот закон, и позволяет увидеть, какой результат дают ограничения на WIP. 

На данный момент есть онлайн-версия игры Featureban\footnote{Онлайн версия игры доступна по ссылке: \url{https://tools.flowfast.io/featureban}, дата обр. 13.01.2026}, но эта версия игры только для одной команды и рассчитана на 3–5 человек. Версия игры на несколько команд существует только в оффлайн формате и разработана сертифицированными экспертами из России (Neogenda)\cite{multi-featureban}. Исходная версия Featureban от Майка Барроуза доступна под лицензией Creative Commons Attribution ShareAlike 4.0~\cite{license-cc}. Все материалы разрабатываемой мультикомандной модификации наследуют этот же правовой статус. Сейчас для проведения игры с несколькими командами приходится адаптировать сторонние платформы, например, Miro или Trello. Однако данные решения не могут быть напрямую использованы для реализации симуляции. 

Во-первых, они только визуализируют игру, но не накладывают никаких ограничений на действия участников. Из-за этого пропадает поддержка ролей, разделение на команды, разрушаются четкие сценарии игр, что являются критически важными аспектами симуляции для корректного моделирования принципов Kanban. Исследования игровых симуляций подчеркивают, что строгость правил, ограничений и сценариев является критически важным аспектом подобных обучающих методов. 

Во-вторых, подобные платформы не собирают нужную статистику. При отсутствии автоматизации сбора статистики увеличивается влияние человеческого фактора, становится сложнее проводить рефлексию над экспериментом, что приводит к снижению педагогической ценности проведения этой игры. 

Таким образом, такие универсальные инструменты не предоставляют достаточно функциональный интерфейс для высокоэффективного проведения мультикомандной версии игры.